\documentclass[11pt,a4paper]{article}

% --- Mathematics (must precede unicode-math) ---
\usepackage{amsmath}
\usepackage{mathtools}

% --- Fonts (lualatex) ---
\usepackage{fontspec}
\usepackage{unicode-math}
\setmathfont{KpMath-Regular.otf}
\setmathfont{KpMath-Bold.otf}[version=bold]

% --- Microtypography (after all font setup) ---
\usepackage{microtype}

% --- Colors & page layout ---
\usepackage[dvipsnames]{xcolor}
\usepackage[margin=25mm, top=30mm, bottom=30mm]{geometry}

% --- Headers & footers ---
\usepackage{fancyhdr}
\pagestyle{fancy}
\fancyhf{}
\fancyhead[L]{\small\itshape P--SV--SH Reciprocity}
\fancyhead[R]{\small\itshape Nestor (2026)}
\fancyfoot[C]{\thepage}
\renewcommand{\headrulewidth}{0.4pt}
\renewcommand{\footrulewidth}{0pt}

% --- Section numbering ---
\setcounter{secnumdepth}{3}

% --- Tables ---
\usepackage{booktabs}
\usepackage{array}

% --- Captions ---
\usepackage[font=small, labelfont=bf, labelsep=period]{caption}

% --- Spacing ---
\usepackage{setspace}
\setstretch{1.15}
\setlength{\parskip}{0.4em}

% --- Bibliography ---
\usepackage[round,authoryear]{natbib}
\bibliographystyle{plainnat}

% --- Hyperlinks (last before macros) ---
\usepackage[colorlinks=true, linkcolor=blue!60!black,
            citecolor=blue!60!black, urlcolor=blue!60!black]{hyperref}

% --- Macros (house set) ---
\newcommand{\bG}{\mathbf{G}}
\newcommand{\bI}{\mathbf{I}}
\newcommand{\bx}{\mathbf{x}}
\newcommand{\dd}{\mathrm{d}}
\newcommand{\bT}{\mathbf{T}}

% --- New macros for this document ---
\newcommand{\bu}{\mathbf{u}}
\newcommand{\bt}{\mathbf{t}}
\newcommand{\bfv}{\mathbf{f}}
\newcommand{\bsig}{\boldsymbol{\sigma}}
\newcommand{\bL}{\mathbf{L}}
\newcommand{\bM}{\mathbf{M}}
\newcommand{\bN}{\mathbf{N}}
\newcommand{\xhat}{\hat{\mathbf{x}}}
\newcommand{\dhat}{\hat{\mathbf{d}}}
\newcommand{\phat}{\hat{\mathbf{p}}}
\newcommand{\qhat}{\hat{\mathbf{q}}}
\newcommand{\nhat}{\hat{\mathbf{n}}}
\newcommand{\kin}{k_{\mathrm{in}}}
\newcommand{\kout}{k_{\mathrm{out}}}
\newcommand{\fout}{\bfv_{\mathrm{out}\leftarrow\mathrm{in}}}
\newcommand{\fin}{\bfv_{\mathrm{in}\leftarrow\mathrm{out}}}

\begin{document}

% --- Title block ---
\thispagestyle{plain}
\begin{center}
\vspace*{1em}
{\LARGE\bfseries P--SV--SH Reciprocity:\\[0.3em]
Mode Conversion, Flux Normalisation and the T-Matrix Metric}\\[1.8em]
{\large T.\,M.\ Nestor}\\[0.5em]
{\itshape Research School of Earth Sciences, Australian National University}\\[2em]
\end{center}

\noindent\textbf{Status.} Explanatory note. The far-field relation
\eqref{eq:ff-recip} is validated numerically for the P$\leftrightarrow$SV
channel of the exact Mie sphere (Section~\ref{sec:mie-check}; Python only,
no Mathematica cross-check yet). The flux normalisation
\eqref{eq:flux-norm} and the T-matrix metric statement
\eqref{eq:t-recip} remain \emph{unvalidated} in this repository.

% =============================================================================
\section{Origin: Betti--Rayleigh reciprocity}
\label{sec:origin}
% =============================================================================

Let the medium be linear elastic with the major symmetry
$c_{ijkl} = c_{klij}$. It may be anisotropic, heterogeneous or
attenuating \citep{Betti1872,Rayleigh1873,Rayleigh1894,Achenbach2003}.
For two time-harmonic states $A$ and $B$ in the same
medium, with no sources inside a closed surface $S$,
\begin{equation}
  \oint_S \bigl( \bu^A \cdot \bt^B - \bu^B \cdot \bt^A \bigr)\,\dd S = 0,
  \qquad \bt = \bsig \cdot \nhat .
  \label{eq:betti}
\end{equation}
Two properties of \eqref{eq:betti} matter below.
\begin{itemize}
\item It is \textbf{bilinear}: there is no complex conjugate. It is
  not an energy statement, so it survives attenuation.
\item It is \textbf{antisymmetric} in $A \leftrightarrow B$, so it is a
  symplectic form on displacement--traction pairs $(\bu, \bt)$.
\end{itemize}
The point-source form is the Green's-tensor symmetry
\begin{equation}
  G_{ij}(\bx, \bx') = G_{ji}(\bx', \bx),
  \label{eq:green-sym}
\end{equation}
which swaps source with receiver and force component with receiver
component \citep{KnopoffGangi1959}; see \citet[\S2.3]{AkiRichards2002}.

% =============================================================================
\section{Mode conversion breaks naive symmetry}
\label{sec:naive}
% =============================================================================

For scalar waves the far-field amplitude obeys
$f(\xhat; \dhat) = f(-\dhat; -\xhat)$. Elastic waves have two speeds,
$\alpha$ (P) and $\beta$ (S), and the converted amplitudes are
\emph{not} symmetric in displacement units.

Write the scattered far field for a unit-amplitude incident plane wave
of direction $\dhat$ and polarisation $\phat$ as
$\fout\, e^{i\kout r}/r$. The far-field Green's tensor carries a
mode-dependent factor $1/(4\pi\rho c^2)$ \citep[\S4.2]{AkiRichards2002}:
\begin{equation}
  \bG^{P} \sim \frac{\xhat\xhat}{4\pi\rho\alpha^2}\,\frac{e^{ik_P r}}{r},
  \qquad
  \bG^{S} \sim \frac{\bI - \xhat\xhat}{4\pi\rho\beta^2}\,\frac{e^{ik_S r}}{r}.
  \label{eq:green-ff}
\end{equation}
Building each plane wave as the limit of a distant point force and
applying \eqref{eq:green-sym} to the scattered part gives
\begin{equation}
  \kin^{2}\; \qhat \cdot \fout(\xhat; \dhat, \phat)
  \;=\;
  \kout^{2}\; \phat \cdot \fin(-\dhat; -\xhat, \qhat).
  \label{eq:ff-recip}
\end{equation}
This is the elastic form of the reciprocity relation for scattering
amplitudes \citep{PaoVaratharajulu1976,Varatharajulu1977}.
Consequently:
\begin{itemize}
\item \textbf{P$\to$P and S$\to$S (including SV$\leftrightarrow$SH).}
  Equal speeds on both sides, so \eqref{eq:ff-recip} is plain symmetry.
  The only subtlety is the sign convention of the SV unit vector under
  direction reversal.
\item \textbf{P$\to$S versus S$\to$P.} These differ by
  $(k_S/k_P)^2 = (\alpha/\beta)^2$, which is $3$ for a Poisson solid.
\end{itemize}

\subsection{Check against the exact Mie sphere}
\label{sec:mie-check}

For a sphere, rotational symmetry reduces \eqref{eq:ff-recip} to a
statement about a single pair of Mie amplitudes. Take P incident along
$+\hat{\mathbf{z}}$ and observe SV at polar angle $\theta$ in the
$xz$-plane, projected on $\hat{\boldsymbol{\theta}}$; call this
$f_{P\to SV}(\theta)$. The reverse experiment has S incident along
$-\xhat(\theta)$ with polarisation $\hat{\boldsymbol{\theta}}$, and P observed
along $-\hat{\mathbf{z}}$. A rotation about $\hat{\mathbf{y}}$ by
$\pi-\theta$ carries it onto SV incident along $+\hat{\mathbf{z}}$ with
polarisation $-\hat{\mathbf{x}}$, observed at $(\theta, \phi=\pi)$. The
$\cos\phi$ dependence of the radial P field cancels the polarisation
sign, and projecting the outgoing P wave along $-\hat{\mathbf{z}}$ onto
$\hat{\mathbf{z}}$ leaves one minus sign:
\begin{equation}
  k_P^{2}\, f_{P\to SV}(\theta) \;=\; -\,k_S^{2}\, f_{SV\to P}(\theta),
  \label{eq:mie-recip}
\end{equation}
where $f_{SV\to P}(\theta)$ is the radial P amplitude for an
$\hat{\mathbf{x}}$-polarised SV wave incident along $+\hat{\mathbf{z}}$.

The gate
\begin{center}
\texttt{scripts/gate\_mie\_psv\_reciprocity.py}
\end{center}
evaluates both sides using \texttt{compute\_elastic\_mie} and
\texttt{mie\_far\_field}.
The P- and SV-incident branches use different right-hand sides and
different far-field assembly (the $m=0$ Legendre expansion versus the
renormalised $m=1$ N-type expansion), so a normalisation fault in either
branch would break \eqref{eq:mie-recip}. The background is
$\alpha = 5$~km/s, $\beta = 3$~km/s, $\rho = 2.5$~g/cm$^3$, with
$(\alpha/\beta)^2 = 2.78$. The gate covers $k_S a \in \{0.1, 0.5, 1.5,
4.0\}$, 60 angles in $[0.05, \pi-0.05]$, and three contrasts:
\emph{moderate} ($\Delta\lambda = +2$~GPa, $\Delta\mu = +1$~GPa,
$\Delta\rho = +100$~kg/m$^3$), \emph{density only}
($\Delta\rho = +250$~kg/m$^3$) and \emph{strong soft}
($\Delta\lambda = -8$~GPa, $\Delta\mu = -10$~GPa,
$\Delta\rho = -500$~kg/m$^3$).

\begin{center}
\begin{tabular}{lcccc}
\toprule
Contrast & $k_S a = 0.1$ & $0.5$ & $1.5$ & $4.0$ \\
\midrule
Moderate     & $9.3\times10^{-5}$ & $1.1\times10^{-5}$ & $2.8\times10^{-6}$ & $5.1\times10^{-6}$ \\
Density only & $5.9\times10^{-5}$ & $5.1\times10^{-6}$ & $1.9\times10^{-6}$ & $3.2\times10^{-6}$ \\
Strong soft  & $1.1\times10^{-4}$ & $1.4\times10^{-5}$ & $6.1\times10^{-6}$ & $5.6\times10^{-6}$ \\
\bottomrule
\end{tabular}
\end{center}

\noindent The table gives the maximum deviation over angle, normalised by
the peak amplitude. The residual comes from \texttt{mie\_far\_field}
evaluating the field at $r = 10^{6}a$ rather than at infinity, an
$\mathcal{O}(n^2/kr)$ near-field remainder that is largest at low
frequency. Two negative controls fail in every case, as they must. The
unweighted scalar form $f_{P\to SV} = -f_{SV\to P}$ is off by exactly
$(\alpha/\beta)^2 - 1 = 1.78$, and reversing the sign of
\eqref{eq:mie-recip} gives an error of $2$.

% =============================================================================
\section{Flux normalisation restores symmetry}
\label{sec:flux}
% =============================================================================

Normalise each amplitude so that its squared modulus is an energy flux.
The incoming and outgoing spherical-wave amplitudes carry
$\sqrt{\rho c}$, and the plane-wave expansion contributes $1/\kin$,
giving
\begin{equation}
  \tilde{\bfv}_{\mathrm{out}\leftarrow\mathrm{in}}
  = \frac{\kin^{3/2}}{\kout^{1/2}}\; \fout .
  \label{eq:flux-norm}
\end{equation}
In these units the scattering matrix over
$(\mathrm{P}, \mathrm{SV}, \mathrm{SH}) \times$ directions is
\textbf{symmetric} under direction reversal,
$\tilde{S} = \tilde{S}^{\mathsf T}$, and in the absence of attenuation it
is also unitary \citep{VaratharajuluPao1976,Varatharajulu1977}.
Reciprocity supplies the symmetry and energy
conservation the unitarity. They are independent facts that share one
normalisation.

The layered-media analogue is given by \citet[Ch.~5]{AkiRichards2002}
and \citet{Kennett1983}: the displacement
coefficients satisfy $R_{PS} \neq R_{SP}$, whereas the
energy-normalised coefficients, with a factor
$\sqrt{\rho\, c \cos\theta}$ per wave, form a symmetric matrix.

% =============================================================================
\section{Decoupling of SH}
\label{sec:sh}
% =============================================================================

In two-dimensional or laterally homogeneous geometries (plane layers,
spheres) SH decouples from P--SV, and reciprocity splits into blocks:
\begin{itemize}
\item \textbf{SH block}: scalar with a single speed, so plainly symmetric;
\item \textbf{P--SV block}: $2\times 2$. The P$\leftrightarrow$SV
  off-diagonal terms need the weighting of \eqref{eq:flux-norm}.
\end{itemize}
For a sphere in the vector spherical-wave basis
\citep{MorseFeshbach1953,BenMenahemSingh1981}:

\begin{center}
\begin{tabular}{llll}
\toprule
Basis field & Character & Wave type & Coupling \\
\midrule
$\bL$ & longitudinal & P  & couples with $\bN$ \\
$\bN$ & poloidal     & SV & couples with $\bL$ \\
$\bM$ & toroidal     & SH & decoupled \\
\bottomrule
\end{tabular}
\end{center}

\noindent Within each multipole order $n$, $\bL$ and $\bN$ form the
P--SV block and $\bM$ stands alone.

% =============================================================================
\section{T-matrix form and the reciprocity metric}
\label{sec:tmatrix}
% =============================================================================

In a basis of regular and outgoing fields, \eqref{eq:betti} becomes the
matrix statement \citep{Waterman1976,Pao1978}
\begin{equation}
  \bT^{\mathsf T} W = W \bT ,
  \label{eq:t-recip}
\end{equation}
where $W$ is the Gram matrix of the bilinear form
$\oint (\bu^A\cdot\bt^B - \bu^B\cdot\bt^A)\,\dd S$ between regular and
outgoing basis fields. That form is the symplectic $J$ acting on
displacement--traction columns.

\begin{itemize}
\item With un-normalised $\bL/\bM/\bN$ columns, $W$ is diagonal with
  mode-dependent entries ($k_P$ for $\bL$, $k_S$ for $\bM$ and $\bN$,
  together with $\rho$ and $\mu$ factors). $\bT$ is then only
  $W$-symmetric, not symmetric.
\item Rescaling the columns by $W^{1/2}$ yields
  $\tilde{\bT} = W^{1/2} \bT W^{-1/2}$ with
  $\tilde{\bT} = \tilde{\bT}^{\mathsf T}$ exactly: the spherical-wave
  counterpart of \eqref{eq:flux-norm}.
\end{itemize}

This is how two operators built with different column normalisations
(T0 and G0 in the IntraPlane Foldy--Lax work) can each be reciprocal yet
appear to disagree: each is symmetric only in its own metric.
Evaluating $W$ from the explicit displacement--stress arrays with the
plain symplectic $J$ checks both against the same form, so any residual
mismatch is a genuine error rather than a normalisation artefact.

\paragraph{Numerical check.} Build $W$ from the explicit arrays at a
test radius and confirm, separately for the P--SV and SH blocks, that
\begin{equation}
  \frac{\lVert \bT^{\mathsf T} W - W \bT \rVert}{\lVert W \bT \rVert}
  \approx \varepsilon_{\mathrm{mach}} .
  \label{eq:check}
\end{equation}

\bibliography{psvsh_reciprocity}

\end{document}
